% Security applications: one per direction, on the same pair of detectors. The pullback compares the GBDT and MalConv
% file by file (CM, 6.1); the pushforward audits MalConv across a corpus (AU, 6.2).
Deployed malware detection layers several static detectors over different representations of the same file. Microsoft
Defender combines feature-specific models in stacked ensembles alongside deep networks that read full file
contents~\cite{microsoft2019insideout}; the winning defense of the 2020 evasion competition combined EMBER-based GBDTs, a
byte skip-gram model, and signatures, so that an attacker has to evade all of them at once~\cite{quiring2020against};
and whether an evasion crafted for one detector transfers to another depends on the pair~\cite{salman2024robustness}.
The members' explanations live in different spaces, features for one and bytes for the other, so they cannot be laid
side by side, and a raw-byte member's explanations cannot be aggregated across files, since a byte offset means
something different in every file. We use each transfer for one defensive task on such a pair. The
pullback, used locally, puts both detectors' explanations on the bytes of a single file and compares them
(Section~\ref{sec:XML:Comparison}). The pushforward, used globally, puts the raw-byte detector's explanations into named
features shared by every file and audits what it relies on across a corpus (Section~\ref{sec:XML:Audit}).
We study the GBDT and MalConv as stand-ins for the two kinds of member, one over engineered features and one over raw
bytes; both are public with their training data, which lets us measure what a defender would see. The method needs
only one explainable model of each kind, not this particular pair.

\paragraph{Setting.} The \emph{defender} operates a feature-based and a raw-byte detector over the same files and owns
both, so it can explain each (TreeSHAP and Integrated Gradients here), and it holds a labeled corpus. It instruments
the feature-based detector's extractor once, which gives $P$ for every file, also for the raw-byte detector, which
never computes these features. The \emph{attacker} wants a malicious file that a detector flags to be classified
benign without changing what the program does. It may rewrite any byte the Windows loader ignores, such as the DOS
header apart from its first and last fields, the DOS stub, the slack at the end of a section, and header fields the
loader does not check, or set a header field to another value the loader accepts: the capability of standard
functionality-preserving attacks on PE detectors~\cite{demetrio2021adversarial}. It does not append or inject content,
pack the file, or change its code; such manipulations evade MalConv more easily~\cite{kolosnjaji2018adversarial,
demetrio2021adversarial} but are not what we measure. We do not execute modified files: which changes the loader
accepts follows the format's documentation. We trust the extractor's parse; a file crafted to be parsed differently by
the extractor and by the loader could misplace $P$, and we leave such files out of scope.

\subsection{Comparing Two Detectors File by File}
\label{sec:XML:Comparison}

The pullback puts the explanations of a feature-based and a raw-byte detector on the same bytes, at byte resolution, so
that a defender can compare them one file at a time. The defender asks four questions, two about the bytes both
detectors rely on and two about the files on which they disagree. Which bytes do both detectors rely on (\emph{shared
reliance}), and does an attacker who evades the raw-byte detector also weaken the feature-based one, file by file
(\emph{attack transfer})? When the detectors disagree, is the flagging detector's evidence content that the other does
not read, or content that it reads the other way (\emph{disagreement triage}), and does this split tell which detector
is right (\emph{adjudication})? Disagreements are where the defender learns most: one detector may have been evaded, or
may have a blind spot or a false positive.

The realistic alternative needs no provenance: explain both models directly over windows of bytes, with occlusion as in
Section~\ref{sec:XML:PullbackFaithfulness}.
% TODO: replace occlusion by the best window explainer of Section 5.1 once WP1 has run. The pullback adds the feature-based detector's explanation at byte
resolution, from one feature extraction per file. For the GBDT, much of what follows lies in named header fields, which
an expert could map to their bytes by hand with a PE parser; the pullback performs this mapping automatically, per file,
for every feature. The comparison itself uses only $P$ and the two byte-level explanations; where the bytes lie
(offsets, which features read them) is a description we add for the reader.

The GBDT's byte-level explanation is the pullback of interventional TreeSHAP and MalConv's is Integrated Gradients.
As window-level alternatives without $P$ we use MalConv's exact occlusion of its 512-byte convolution windows and the
GBDT's occlusion from Section~\ref{sec:XML:PullbackFaithfulness}, and both models score every deleted file. Deleting
by MalConv's own Integrated Gradients removes far more evidence than deleting the whole file (19 log-odds at 1\% of the
bytes), the fragility of MalConv noted in Section~\ref{sec:XML:PushforwardFaithfulness}, so every comparison is made
at a fixed budget against controls of the same size.
Disagreements are common: at their 1\% false-positive thresholds, the two detectors disagree on 6.4\% of the test split
(5.7--5.8\% for the other two seeds), including 29{,}267 malware detected only by the GBDT, 1{,}964 detected only by
MalConv, and 1{,}547 false positives of each model that the other does not share. Shared reliance uses the 100 files of
Section~\ref{sec:XML:PullbackFaithfulness}; the triage draws 25 test files from each of the four cells of
disagreement.

\paragraph{Shared reliance.} At first sight the two detectors look at different things: their explanations barely
correlate at any resolution (median Spearman $\rho \leq 0.07$), and their top 1\% of bytes overlap at 6.5 times
chance. Yet the bytes the GBDT relies on carry much of MalConv's evidence, and this is visible only at byte level.
Deleting the GBDT's top 0.1\% of bytes by its pullback removes a median 1.54 of MalConv's log-odds, against 0.55 for
the same pullback summed over 256 windows, 0.79 for the GBDT's own occlusion, and nothing for random windows; the
byte-level pullback stays ahead of both window rankings at every budget up to 5\% ($p \leq 0.025$).

To locate the shared reliance, let $T_g$ and $T_m$ be the top 1\% of each model's deletable bytes, and
$S = T_g \cap T_m$ the bytes both rely on (6.3\% of each top set, against 1\% by chance). Deleting $S$ removes a median
2.78 of the GBDT's log-odds and 3.78 of MalConv's. As many bytes drawn from the GBDT's exclusive top set remove 0.86
and $-0.08$, and from MalConv's, 0.04 and 2.58. Byte for byte, the shared bytes are worth more to \emph{each} model
than its own exclusive bytes: they are where a single manipulation weakens both detectors. They lie in the first
kilobyte (97\% of the GBDT's evidence in $S$ and 75\% of MalConv's), and 86\% of the GBDT's credit on them comes from
header and data-directory fields, such as the DLL characteristic, the sizes of the certificate table and debug
directory, and the checksum, consistent with MalConv's known reliance on the header~\cite{demetrio2019explaining}.

\paragraph{Attack transfer.} Deletion shows where shared reliance lies, not what an attacker gains from it. We
therefore attack MalConv as the setting allows and ask whether the pullback predicts, file by file, how much the attack
also costs the GBDT. The attack is a white-box gradient attack on MalConv's byte
embeddings~\cite{kolosnjaji2018adversarial, kreuk2018deceiving} that may rewrite only the loader-ignored bytes of the
DOS header, the DOS stub, and the slack at the end of sections~\cite{demetrio2019explaining, suciu2019exploring,
demetrio2021adversarial}; it runs for a fixed budget of queries and stops early once MalConv's score falls below its
threshold. Because this editable region $E$ is fixed before the attack, the GBDT's pullback on the original file can
place evidence on it, and its sum over $E$ says how much of the GBDT's verdict rests on bytes this attacker may change,
before any attack is run. We compare it, as a predictor of the drop in the GBDT's log-odds on the adversarial file,
with the same sum from occlusion's windows and with the number of bytes the attack changed.

\begin{table}[t]
  \centering
  \small
  \begin{tabular}{lcc}
    \toprule
    Predictor of the GBDT's drop & Spearman $\rho$ & AUC \\
    \midrule
    GBDT evidence on $E$, pullback & \synthetic{0.58} & \synthetic{0.81} \\
    GBDT evidence on $E$, occlusion & \synthetic{0.31} & \synthetic{0.66} \\
    Bytes changed by the attack & \synthetic{0.12} & \synthetic{0.55} \\
    \bottomrule
  \end{tabular}
  \caption{Attack transfer: an attack on MalConv that rewrites only loader-ignored bytes (the region $E$), and how well
  the GBDT's evidence on $E$ in the original file predicts the drop in the GBDT's log-odds. AUC: for a drop of more than
  one log-odds. \synthetic{Placeholder values.}}
  \label{tab:XML:Transfer}
\end{table}

On \synthetic{100} held-out malware that both detectors detect, \synthetic{the attack evades MalConv on 74\% of them and
takes the GBDT's malware probability below its threshold on 9\%}. The pullback predicts which files these are
(Table~\ref{tab:XML:Transfer}): \synthetic{its evidence on $E$ correlates with the GBDT's drop at Spearman 0.58,
against 0.31 over windows and 0.12 for the number of bytes changed}. The editable region is small and lies next to
content the GBDT reads in different ways, such as the Rich header inside the DOS stub, which EMBER summarizes in named
features, and section data around the slack; windows blur these together. A defender can thus tell, per file, which
of the feature-based detector's detections an attacker who evades the raw-byte detector would also put at risk.

\begin{table}[t]
  \centering
  \small
  \begin{tabular}{llccc}
    \toprule
    Flagged by & Label & blind & oppose & agree \\
    \midrule
    MalConv only & malware & 0.77 & 0.10 & 0.13 \\
    MalConv only & benign & 0.77 & 0.10 & 0.14 \\
    GBDT only & malware & 0.25 & 0.33 & 0.42 \\
    GBDT only & benign & 0.46 & 0.17 & 0.37 \\
    \midrule
    \multicolumn{5}{l}{\textit{256 windows, GBDT only, malware:}} \\
    \quad pullback per window & & 0.19 & 0.06 & 0.75 \\
    \quad GBDT occlusion & & 0.28 & 0.02 & 0.70 \\
    \bottomrule
  \end{tabular}
  \caption{Disagreement triage (25 files per cell): the flagging model's evidence in its top 1\% of bytes, split by how
  the other model reads it (mean shares). Bottom: the same split over 256 windows, with the GBDT's window evidence from
  its pullback summed per window or from its occlusion (no $P$). At windows, the oppose share nearly vanishes.}
  \label{tab:XML:Triage}
\end{table}

\paragraph{Disagreement triage.} Shared reliance concerns files on which both detectors agree; triage concerns those on
which they do not. Let A be the model that flags a file and B the other. We take A's top 1\% of deletable bytes by
evidence for malware and split them by how B reads each byte: \emph{blind} if it lies outside B's attended bytes (its
top 5\% by absolute evidence), \emph{oppose} if B attends to it as evidence against malware, and \emph{agree} if as
evidence for it. If the split is right, deleting the blind bytes lowers A's score and leaves B's unchanged, while
deleting the oppose bytes raises B's, since they carry B's verdict.

Table~\ref{tab:XML:Triage} shows that the two directions of disagreement differ in kind. When MalConv alone flags a
file, whether malware or a false positive, about three quarters of its evidence lies on bytes the GBDT does not read,
mostly in the body of the file (57--73\% of it beyond the first 64\,KB): a gap in what the GBDT's features cover. When
the GBDT alone flags malware, MalConv reads most of the same bytes, in the header (87--97\% of the oppose and agree
evidence lies in the first kilobyte), and reads a third of the GBDT's evidence the other way: a disagreement about what
the same content means. Deleting these oppose bytes makes MalConv flag 19 of the 25 files, against 1 of 25 for a random
part of the GBDT's top bytes of the same size (17 against 1 for the GBDT's false positives, and similar with random
fill), while deleting the blind bytes leaves MalConv unchanged. Two cells are descriptive only: where A is MalConv, any
large deletion raises the GBDT's score, which confounds B's response, and the GBDT's credit on its blind bytes sits
mostly on aggregate features (counts of string patterns, hashed section entropies and imports), which deleting a few of
their bytes barely changes.

Each label points the defender to a different response. A disagreement whose evidence is mostly \emph{blind}, as when
MalConv alone flags a file, is a coverage gap: the flagging detector reads content the other's features never see, the
remedy is a feature over that content, and the blind bytes show where it lies. A disagreement whose evidence is mostly
\emph{oppose}, as when the GBDT alone flags malware, is a conflict over the same bytes: the oppose bytes hold the other
detector's benign verdict, so they are the bytes to inspect first. The labels describe which content drives a disagreement; why it
arose is for the analyst who inspects that content.

This diagnosis depends on byte resolution. Repeating the split and the deletions over 256 windows for the 50
GBDT-flagged files, with the GBDT's window evidence taken either from its pullback summed per window or from its
occlusion (Table~\ref{tab:XML:Triage}, bottom), the oppose category nearly vanishes: for GBDT-only malware it falls from
0.33 of the GBDT's top evidence to 0.06 and 0.02, it is nonempty in 6 and 2 of the 25 files instead of all 25, and
deleting it no longer makes MalConv flag more often than a random part of the top set (3 of 6, and 0 of 2). The agree
share absorbs it. Because the oppose and agree bytes interleave within the header, any window merges them, and the
window-level diagnosis reads the opposite of the byte-level one: MalConv appears to agree with the GBDT's evidence.
Since the coarsened pullback behaves like occlusion, the loss comes from resolution, which for the feature-based
detector only the pullback affords.

\paragraph{Adjudication.} The split also helps decide a disagreement. When the GBDT alone flags a file, the split of its
evidence differs with the truth: in the files of Table~\ref{tab:XML:Triage}, the oppose share is 0.33 for malware and
0.17 for false positives, and the blind share 0.25 against 0.46. We test whether this tells which detector is right
beyond what the detectors' own scores say. On \synthetic{200} further GBDT-only disagreements, half malware and half
benign, a logistic regression predicts whether the GBDT is right from the two detectors' log-odds alone, or from the
log-odds and the blind and oppose shares, under repeated five-fold cross-validation. We report the share of the GBDT's
false positives that it removes while keeping 99\% of the GBDT-only malware, at the test split's base rate of 29{,}267
GBDT-only malware to 1{,}547 GBDT-only false positives. \synthetic{The scores alone remove 12\% of the false positives
(AUC 0.71); with the byte-level triage, 31\% (AUC 0.83); with the same triage over windows, 14\% (AUC 0.73).} When
MalConv alone flags a file, the split is the same for malware and benign files (blind 0.77 in both), so it cannot
decide those disagreements, which are coverage gaps either way.

\subsection{Auditing a Raw-Byte Detector}
\label{sec:XML:Audit}

A raw-byte detector such as MalConv can be explained file by file with Integrated Gradients, but these explanations
have no vocabulary in which its reliance recurs across files. Auditing it
for reliance on content that an attacker could change cheaply, without affecting what the program does, is therefore a
corpus-level question~\cite{arp2022and}, and per-file byte explanations cannot be aggregated to answer it. The
pushforward supplies the shared vocabulary. It sends each file's byte explanation into the named features of a feature
extractor, which are the same for every file, splitting the credit of shared bytes among the features that read them.

\paragraph{Threat model.} The attacker is that of the setting. It needs neither MalConv's parameters nor its
explanations, only to know which content MalConv relies on, which an audit like ours reveals. The defender's goal is to
find that content first: what MalConv's detections rest on and an attacker could change cheaply.

The realistic alternative is the surrogate of Section~\ref{sec:XML:PushforwardFaithfulness}, a GBDT fit to MalConv's
log-odds on the EMBER features, explained with TreeSHAP and aggregated over the corpus, with no provenance. It finds
what \emph{predicts} MalConv, not what MalConv \emph{reads}: per file, its explanation carries half as much of MalConv's
evidence as the pushforward's. Its training pool may also
contain relatives of the explanation files, which can only favor it. A hand-built aggregation, a PE parser locating
every field in every file and summing Integrated Gradients over it, would answer the same question, but only for one
file format and at great manual effort; the pushforward does it automatically for any feature extractor, and needs no
second model.

We explain MalConv with Integrated Gradients on the 2{,}000-file explanation set and push each file's explanation
forward through that file's $P$. In each file we take the 20 features with the most credit for malware; across files, a
feature scores the number of files it appears in times its mean credit. The 20 highest-scoring features whose credited
content in a file can be located, at most 256 runs of bytes, are the candidates; this excludes aggregates such as the
byte histogram, but not features whose content is spread over the file or shared with other features, such as a
section's name, which several section features read, or an imported function's name, which import and string
features share. A
candidate is \emph{confirmed} by scrambling: in 191 held-out malware, we overwrite the bytes that the file's own $P$
assigns to the feature with random bytes and compare the drop in MalConv's malware log-odds, paired per file, with that
from scrambling, for every byte scrambled, a random byte elsewhere in the same convolution window. A feature is confirmed if the median drop
and the median paired difference both exceed $0.1$ log-odds (Wilcoxon $p < 0.05/20$). The surrogate's candidates are
found and confirmed in the same way.

\begin{table}[t]
  \centering
  \small
  \begin{tabular}{@{}llcc@{}}
    \toprule
    Field & Loader & Effect & Surrogate \\
    \midrule
    CheckSum & ignored & 0.24 & \checkmark \\
    SizeOfHeapReserve & valid value & 0.16 & \\
    ImageBase & valid value & 0.15 & \checkmark \\
    Subsystem & valid value & 0.14 & \checkmark \\
    Debug directory address & ignored & 0.13 & \checkmark \\
    Global-pointer directory address & ignored & 0.12 & \\
    Export directory address & ignored\textsuperscript{a} & 0.12 & \\
    Copyright directory address & ignored & 0.11 & \\
    Import address table address & needed & 0.10 & \\
    \bottomrule
  \end{tabular}
  \caption{The 9 fields that MalConv relies on, found through the pushforward and confirmed by scrambling, all outside
  the DOS header. Loader: how the Windows loader treats the field. Effect: the median extra drop in MalConv's malware
  log-odds from scrambling the field rather than as many other bytes of the same window. Surrogate: also confirmed from
  the surrogate. \textsuperscript{a}In executables, not in DLLs.}
  \label{tab:XML:AuditFields}
\end{table}

Of MalConv's 20 candidates, 9 are confirmed (Table~\ref{tab:XML:AuditFields}): the checksum, the image base, the size of
the heap reserve, the subsystem, and the addresses of five data directories. All lie outside the
DOS header, the region whose manipulation MalConv is known to be sensitive to~\cite{demetrio2019explaining}: the audit
finds reliance beyond what was documented. The surrogate's candidates recover 4 of these 9 fields and add 5 others
(certificate-table entries, the signing time, the DLL characteristic), with some larger effects. The two methods thus
find reliance of comparable size but only partly the same: each confirms fields the other misses.
\synthetic{Beyond header fields, the audit confirms reliance on two section names and three hashed import buckets,
whose bytes other features read as well; there the pushforward splits each byte's credit among the features that read
it, which summing attributions over parsed regions does not do.}
% TODO: Table and counts above are from the run restricted to contiguous content (runs 44-49); update after WP7.

Scrambling measures reliance; whether that reliance can be exploited is a separate question, which we answer with
changes an attacker could actually make. In 1{,}000 fresh held-out malware, none in the explanation set or used for
confirmation, of which MalConv detects 902 at its 1\% false-positive threshold, we change the confirmed fields in
loader-safe ways only (Table~\ref{tab:XML:AuditFields}). Fields the loader ignores are scrambled; fields it checks
take a valid value drawn from benign files of the same machine type (an image base kept 64K-aligned and changed only
when the file has relocations, a GUI or console subsystem, a heap reserve consistent with its commit); the import
address table's address, which the loader needs, is left alone. The files are not executed; the classification of
fields follows the format's documentation. On average 7.3 of the 9 fields can be changed this way in a file, and doing
so evades MalConv on 12.4\% of the detected malware (95\% CI 10.3--14.6). That is about twice as often as changing as
many other loader-ignored fields outside the DOS header, drawn per file from the 74 that no method confirmed (6.0\%;
paired difference $+6.4$ points, $p = 3 \cdot 10^{-12}$). Under the same rules, the surrogate's fields
evade on 10.4\% ($p = 0.027$ against the pushforward's), and the union of the two methods' fields on 18.4\%. The
pushforward is thus the more faithful and the stronger single method, needs no second model, and complements the
surrogate: a defender auditing MalConv should use both.

The audit also tells the defender what to change: a field that MalConv relies on but an attacker can change freely
should not be read. We test this at test time, without retraining. Before MalConv scores a file, the audited fields an
attacker can change (all confirmed fields but the import address table's, which the loader needs) are overwritten with
zeros at the bytes the file's own $P$ names; MalConv's threshold is re-set to a 1\% false-positive rate on 5{,}000
held-out benign files, and the attacks above are repeated on the malware each defense still detects
(Table~\ref{tab:XML:Mitigation}). An attack on a field the defense zeroes fails by construction, so the informative
quantities are the cost in clean detection and the strongest attack that remains. Zeroing the pushforward's 8 fields
lowers detection from 91.2\% to 89.2\% and cuts the strongest loader-safe attack from 16.8\% to 8.2\%; zeroing the 12
fields of both methods costs 3.5 points and cuts it to 2.9\%. The fields thus carry real signal as well as an attack
surface, and removing them is a trade-off rather than a free fix. Defending only part of the reliance also shifts the
attack: with the pushforward's fields zeroed, changing other loader-ignored fields evades more often (8.2\% against
5.6\%), not because each change moves MalConv further but because the detected malware now sit closer to the re-set
threshold. What the audit provides is the list of fields to weigh, found automatically; the defense that uses both
methods' fields removes most of the attack surface we can measure. \synthetic{Retraining MalConv on files with the
fields zeroed lowers evasion to $1.8\%$ with no measurable clean-detection cost.}

\begin{table}[t]
  \centering
  \small
  \begin{tabular}{lccc}
    \toprule
    MalConv input & detection & strongest & other fields \\
    \midrule
    unmodified & 91.2\% & 16.8\% & 5.6\% \\
    zero pushforward fields & 89.2\% & 8.2\% & 8.2\% \\
    zero both methods' fields & 87.7\% & 2.9\% & 2.9\% \\
    retrain, fields zeroed & \synthetic{90.9\%} & \synthetic{1.8\%} & \synthetic{--} \\
    \bottomrule
  \end{tabular}
  \caption{The audit's countermeasure: the audited fields an attacker can change are zeroed before MalConv scores a
  file. Detection of 1{,}000 held-out malware at a 1\% false-positive rate, re-set per row on 5{,}000 benign files;
  evasion of the detected malware by the strongest loader-safe attack, and by other loader-ignored fields alone.
  \synthetic{The retrained row is a placeholder.}}
  \label{tab:XML:Mitigation}
\end{table}

The two applications use the transfers in opposite directions on the same pair of detectors. File by file, the
pullback shows which bytes both detectors rely on, predicts which detections an attacker who evades the raw-byte
detector also puts at risk, and helps decide which detector is right when they disagree; across the corpus, the pushforward names what the raw-byte detector relies on, and loader-safe changes to it evade the
detector. Neither needs a hand-built mapping from features to bytes.
